\documentclass[11pt,a4paper]{article}
\usepackage[T1]{fontenc}
\usepackage[utf8]{inputenc}
\usepackage{amsmath,amssymb,xcolor}
\usepackage[margin=25mm]{geometry}
\usepackage[colorlinks=true,urlcolor=blue]{hyperref}
\definecolor{nexorablue}{HTML}{2459B5}
\setlength{\parindent}{0pt}
\setlength{\parskip}{7pt}
\setlength{\emergencystretch}{2em}
\renewcommand{\familydefault}{\sfdefault}
\begin{document}

{\small\bfseries\color{nexorablue} NEXORA CPGE}\par
{\small MP / NEX-IND-03 / Énoncé}\par
{\LARGE\bfseries Tension d’une bobine en rotation}\par
\vspace{0.25em}\hrule\vspace{0.6em}

Une bobine de $40$ spires d’aire $0{,}012\,\mathrm{m^2}$ tourne à $30\,\mathrm{rad\,s^{-1}}$ dans un champ de $0{,}15\,\mathrm T$. La normale coïncide avec le champ à $t=0$. Donner la tension induite orientée, son amplitude et sa valeur efficace.


\end{document}
